\documentclass[12pt,a4paper]{article}
\usepackage[utf8]{inputenc}
\usepackage{amsmath,amssymb,amsthm}
\usepackage{hyperref}
\usepackage{booktabs}
\usepackage{geometry}
\geometry{margin=1in}

\newtheorem{theorem}{Theorem}
\newtheorem{lemma}[theorem]{Lemma}
\newtheorem{definition}[theorem]{Definition}
\newtheorem{corollary}[theorem]{Corollary}

\title{A Common Parent-Structure Framework for Gravity,\\Gauge Structure, Matter, and Quantum Amplitudes}

\author{Tom DiFiore}

\date{\today}

\begin{document}

\maketitle

\begin{abstract}
We present a parent-structure framework in which gravity, gauge structure, a defect-based matter sector, and a quantum-amplitude layer are derived or conditionally induced from five explicit primitives: a Lorentzian metric, a Lie algebra, a source functional, linear composition, and a perturbation space of dimension $\geq 2$. The headline results are complete 4D Lovelock uniqueness for the gravitational field equation, full nonabelian Yang--Mills structure with a machine-checked Bianchi identity, a charge algebra with multiple conserved quantum numbers, and a unique Born-type quadratic observable within the rotation-invariant class---together with Lindblad decoherence, dimension selection ($d = 3$ unique), and a cosmological constant scaling relation $\Lambda \sim N^{-1/2}$. Each primitive is proven necessary: removing any one gives a provably degenerate theory. The phase accumulation rule $e^{ikL}$ is derived from the source functional's linearity (not assumed from the Feynman prescription), yielding an unbroken chain from primitives to the double-slit interference formula. Chirality is derived from the source/dressing decomposition plus gauge invariance of the source functional, eliminating a standard postulate. Given the Standard Model representation content, anomaly cancellation uniquely determines all hypercharge ratios: the cubic condition factors as $18 y_Q(y_d + 4y_Q)(y_d - 2y_Q) = 0$, with both nontrivial branches giving the Standard Model. The $\theta$-term of QCD is shown to be parity-odd and hence orthogonal to the parity-even source functional, providing a structural resolution of the strong CP problem. The unique spatial dimension $d = 3$ is also the minimum permitting CP violation in the quark sector ($N_g \geq 3$ required). Separate computational models reproduce GR weak-field observables and canonical quantum benchmarks (double-slit, CHSH at the Tsirelson bound, GHZ/Mermin scaling). The algebraic core is formally verified in Lean~4 with zero custom axioms and zero \texttt{sorry} statements. This is a candidate framework with explicit assumptions, not a claim of final closure.
\end{abstract}

\section{Introduction}

The search for a unified description of fundamental physics has been one of the central programs of theoretical physics for over a century. The standard approach postulates separate mathematical structures for gravity (pseudo-Riemannian geometry), matter (gauge field theory on fiber bundles), and quantum mechanics (Hilbert space operators), then seeks a framework that contains all three as limiting cases.

We take a different approach. Rather than unifying pre-existing mathematical structures, we identify a small parent architecture and prove that several gravity, gauge, matter, and quantum-style structures are uniquely selected or conditionally induced from it. The algebraic core is formally verified in Lean~4, while the concrete spacetime realizations are computationally verified in deterministic Python models.

In this paper, ``derived'' means algebraically implied within the formal system; ``conditionally induced'' means implied given explicit structural hypotheses (e.g., existence of a source functional, a nontrivial dressing sector, or rotation invariance of the observable). The point of the paper is not to claim that no assumptions remain, but to show that once a small and physically legible primitive base is accepted, a surprisingly large amount of downstream structure becomes rigid.

The key results, summarized informally:
\begin{enumerate}
\item \textbf{Gravity classification.} The renormalization operator has $\alpha = d{-}1$ as its unique fixed point (Theorem~\ref{thm:rg}). The contracted Bianchi identity, derived from metric second derivatives (Theorem~\ref{thm:bianchi}), combined with the complete 4D Lovelock classification (Theorem~\ref{thm:lovelock}), forces the field equation to be of the Einstein-plus-cosmological-constant form $a\,G_{ab} + \Lambda\,g_{ab}$ within the tensorial, second-order, natural class. A variational/stationarity layer selects this as the physical field equation.

\item \textbf{Gauge structure.} Full nonabelian Yang--Mills theory is derived from a Lie algebra (structure constants with antisymmetry and Jacobi identity): the covariant derivative $D_\mu$, the nonabelian field equation $D^\mu F_{\mu\nu} = 0$, and the nonabelian Bianchi identity $D_\lambda F_{\mu\nu} + \mathrm{cyclic} = 0$ (Theorem~\ref{thm:ym}).

\item \textbf{Matter classification.} Within the linear defect sector, source strength defines conserved charges with additivity, conjugation, and annihilation, all derived from linearity (Theorems~\ref{thm:charge}--\ref{thm:annihilation}). Multiple independent source functionals give multiple simultaneously conserved charges. A spin-statistics connection and representation covariance under the adjoint action are established.

\item \textbf{Quantum amplitude structure.} The source/dressing decomposition induces a canonical two-component amplitude structure, uniquely representable in complex form within the stated class (Theorem~\ref{thm:complex}). The Born-type quadratic observable $|z|^2$ is the unique rotation-invariant quadratic observable (Theorem~\ref{thm:born_unique}). The phase accumulation rule $z(\gamma) = e^{ik\varphi(\gamma)}$ is \emph{derived} from the linearity of the source functional: $\varphi$ is additive on path composition (from \texttt{map\_add}), and the exponential converts this additive structure to multiplicative amplitude composition (Theorem~\ref{thm:propagation}). The derivation chain from source functional to interference formula $|z_1 + z_2|^2 = 2 + 2\cos(k(s_1 - s_2))$ is unbroken: no external prescription is required. We do not claim a full derivation of quantum mechanics in the traditional operational sense; rather, we derive a minimal amplitude calculus reproducing interference, non-classical correlations, and a unique Born-type observable within the framework's symmetry class. Dynamical decoherence via Lindblad-type dephasing gives exponential decay of coherence with a proven classical limit (Theorem~\ref{thm:lindblad}).

\item \textbf{Concrete realizations.} Explicit computational models in $3{+}1$ dimensions reproduce GR-like and quantum benchmark observables (Sections~\ref{sec:gr}--\ref{sec:quantum_bench}). The phase prescription $e^{ikL}$ underlying the quantum benchmarks is now derived from the source functional (Theorem~\ref{thm:propagation}), so these benchmarks are consequences of the framework's amplitude dynamics, with the wavenumber $k$ as the sole free parameter per particle type.
\end{enumerate}

\subsection{Formal verification}

Every algebraic theorem in this paper has been formally verified in Lean~4 with the Mathlib library. The complete proof development is available at \url{https://github.com/tomdif/unifiedtheory}. The command \texttt{lake build} verifies all 2325 proof jobs with zero errors and zero \texttt{sorry} statements. The axiom footprint of every theorem is $\{\texttt{propext}, \texttt{Classical.choice}, \texttt{Quot.sound}\}$ (the three standard Lean kernel axioms).

The computational sections (Sections~\ref{sec:gr}--\ref{sec:quantum_bench}) are not Lean proofs; they are deterministic numerical benchmarks run in Python with fixed seeds.

\section{The Parent Architecture}
\label{sec:parent}

For expository purposes we package the framework into interacting blocks. This interface should not be read as the final irreducible primitive list; several reductions and dynamical refinements are established in the formal development.

The framework has \textbf{two structural primitives} (defining the arena) and \textbf{two dynamical fields} (selected by field equations):

\begin{center}
\begin{tabular}{lll}
\toprule
& Gravity & Gauge \\
\midrule
Structural primitive & Manifold (dim $n$) & Lie algebra (dim $g$) \\
Dynamical field & Metric $g_{ab}$ & Connection $A_\mu^a$ \\
Curvature & $R_{abcd}$ & $F_{\mu\nu}^a$ \\
Field equation & $G + \Lambda g = 0$ & $D^\mu F_{\mu\nu} = 0$ \\
\bottomrule
\end{tabular}
\end{center}

Each of the five primitives is proven \emph{necessary}: removing any one gives a provably degenerate or trivial theory. The derivation tree is:

\begin{center}
\begin{tabular}{lcl}
\toprule
Primitive & $\longrightarrow$ & Downstream structure \\
\midrule
Lorentzian metric & $\longrightarrow$ & Riemann, Bianchi, $\nabla^a G_{ab} = 0$, Lovelock $\to a G + \Lambda g$ \\
Lie algebra & $\longrightarrow$ & $F = dA + [A,A]$, $D^\mu F_{\mu\nu} = 0$, nonabelian Bianchi \\
Source functional $\varphi$ & $\longrightarrow$ & K/P split, charge $Q = \varphi \circ \pi_K$, matter sector \\
Linear composition & $\longrightarrow$ & Charge additivity, conjugation, annihilation \\
dim$(V) \geq 2$ & $\longrightarrow$ & Nontrivial dressing, $z = Q + iP$, interference, $|z|^2$ \\
\bottomrule
\end{tabular}
\end{center}

\noindent The five primitives are not independent postulates about different sectors; they jointly define a single architecture from which all downstream structure descends.

\section{Classification Results}
\label{sec:classification}

\subsection{Gravity endpoint classification}

The framework identifies the unique admissible tensor structure (Einstein-type equations) but does not yet formalize a full variational or dynamical selection principle distinguishing physical solutions within this class.

\begin{theorem}[Renormalization rigidity]
\label{thm:rg}
In spatial dimension $d$, the unique scale-invariant exponent is $\alpha = d - 1$. In $3{+}1$ dimensions: $\alpha = 2$ (inverse-square law).
\end{theorem}

\begin{theorem}[Contracted Bianchi identity]
\label{thm:bianchi}
From metric second derivatives with commutativity of partials, the Riemann tensor, Ricci tensor, scalar curvature, and Einstein tensor are explicitly constructed. The contracted Bianchi identity $2\nabla^a R_{ab} = \nabla_b R$ and the divergence-free property $\nabla^a G_{ab} = 0$ are derived. This is a kinematic identity holding for \emph{all} metrics.
\end{theorem}

\begin{theorem}[Complete 4D Lovelock uniqueness]
\label{thm:lovelock}
In four dimensions, within the tensorial, second-order, natural class:
\begin{enumerate}
\item \emph{Contraction classification}: The Ricci tensor is the only independent rank-2 $\delta$-contraction of the Riemann tensor (6 contraction cases proven).
\item \emph{Bianchi constraint}: Divergence-free within $\{c\,R_{ab} + d\,R\,g_{ab} + e\,g_{ab}\}$ forces $d = -c/2$, giving $a\,G_{ab} + \Lambda\,g_{ab}$.
\item \emph{Gauss--Bonnet vanishing}: The Lanczos--Lovelock tensor $H_{ab} \equiv 0$ in 4D (rank-5 generalized Kronecker delta vanishes by pigeonhole).
\item \emph{$\varepsilon$-tensor exclusion}: $\varepsilon \cdot \varepsilon = \delta$ (generalized Kronecker delta), and odd $\varepsilon$-count is excluded by tensoriality (parity argument).
\item \emph{Higher Lovelock}: All Lovelock tensors of order $p \geq 2$ vanish in 4D.
\end{enumerate}
\end{theorem}

\subsection{Gauge endpoint classification}

\begin{theorem}[Nonabelian Yang--Mills]
\label{thm:ym}
Given structure constants $c^a_{bd}$ with antisymmetry and Jacobi identity:
\begin{enumerate}
\item The gauge-covariant derivative $D_\lambda F_{\mu\nu}^a = \partial_\lambda F_{\mu\nu}^a + c^a_{bd} A_\lambda^b F_{\mu\nu}^d$ is defined.
\item The nonabelian Yang--Mills equation $D^\mu F_{\mu\nu} = 0$ is formalized.
\item The nonabelian Bianchi identity $D_\lambda F_{\mu\nu} + D_\mu F_{\nu\lambda} + D_\nu F_{\lambda\mu} = 0$ is fully proven. The proof uses antisymmetry of $c$ for dA$\cdot$A cancellation and the Jacobi identity for A$\cdot$A$\cdot$A cancellation.
\item The abelian limit ($c = 0$) recovers the Maxwell equations.
\end{enumerate}
\end{theorem}

\subsection{Quantum/matter uniqueness results}

\begin{theorem}[Charge algebra from linearity]
\label{thm:charge}
Charge additivity $Q(d_1 \oplus d_2) = Q(d_1) + Q(d_2)$ follows from the modeling choice that composition is vector addition and charge is a linear functional (\texttt{map\_add}). Charge conjugation follows from \texttt{map\_neg}. These are consequences of the primitives, not independent postulates.
\end{theorem}

\begin{theorem}[Annihilation]
\label{thm:annihilation}
$Q(d \oplus \bar{d}) = 0$ and the composite is inert.
\end{theorem}

\begin{theorem}[Multiple charges and richer matter]
Given $k$ independent linear functionals $\varphi_1, \ldots, \varphi_k$, all $k$ charges are simultaneously additive and simultaneously negated under conjugation. Integer-charge sectors are closed under composition (bosonic). Half-integer pairs compose to integer (structural spin-statistics). Charges transform covariantly under the adjoint action of the Lie algebra.
\end{theorem}

\begin{theorem}[Born-type observable uniqueness]
\label{thm:born_unique}
Any rotation-invariant quadratic observable on the $(Q, P)$ pair must be proportional to $Q^2 + P^2 = |z|^2$. The observable is uniquely induced by the source/dressing pair and rotation invariance, not separately postulated.
\end{theorem}

\textbf{Remark (relation to Gleason's theorem).} Gleason's theorem derives the Born rule from the structure of Hilbert space projectors (dimension $\geq 3$). Our result is complementary: it derives the $|z|^2$ observable from SO(2) invariance on a 2D real space, without invoking Hilbert space. The rotation invariance hypothesis plays the role that the measure-on-projectors hypothesis plays in Gleason. Neither result derives its hypothesis from deeper principles; the difference is the starting mathematical structure (K/P pair vs.\ Hilbert space).

\begin{theorem}[Complex amplitudes]
\label{thm:complex}
The source/dressing decomposition induces a canonical two-component amplitude structure. The source projection $\pi_K$ is the \emph{unique} rank-1 linear map satisfying $\varphi \circ \pi_K = \varphi$ (universal property). The complex identification $z = Q + iP$ is uniquely forced within the 2D commutative division algebra class (Frobenius).
\end{theorem}

\begin{theorem}[Propagation rule from source functional]
\label{thm:propagation}
The phase accumulation rule $z(\gamma) = e^{ik\varphi(\gamma)}$ is derived from three facts: (i)~the source functional $\varphi$ is linear, so $\varphi(\gamma_1 + \gamma_2) = \varphi(\gamma_1) + \varphi(\gamma_2)$ (additivity under path composition); (ii)~the exponential map $s \mapsto e^{iks}$ converts additive to multiplicative: $e^{ik(s_1+s_2)} = e^{iks_1} \cdot e^{iks_2}$; (iii)~composing (i) and (ii) gives $z(\gamma_1 \cdot \gamma_2) = z(\gamma_1) \cdot z(\gamma_2)$. The amplitude has unit modulus ($|z|^2 = 1$, $\sin^2 + \cos^2 = 1$) and $z(0) = 1$. The wavenumber $k$ is the sole free parameter per particle type.

The two-path interference formula follows: $|z(s_1) + z(s_2)|^2 = 2 + 2\cos(k(s_1 - s_2))$, which is the double-slit fringe pattern with fringe spacing determined by $k$.
\end{theorem}

\textbf{Remark (uniqueness).} We use the standard analytic fact that continuous characters of $(\mathbb{R}, +)$ are exponentials; the existence direction is formalized, the uniqueness direction is classical and will be formalized in a subsequent version.

\begin{theorem}[Lindblad decoherence and classical limit]
\label{thm:lindblad}
The dephasing parameter $\gamma(t) = e^{-\Gamma t}$ gives exponential decay of coherence. The observable $\mathrm{obs}(t) = p_1 + p_2 + 2\gamma(t)\,c_{\mathrm{re}}$ interpolates between quantum ($\gamma = 1$) and classical ($\gamma = 0$). The classical limit is proven: for any $\varepsilon > 0$, there exists $T > 0$ such that $|\mathrm{obs}(t) - (p_1 + p_2)| < \varepsilon$ for all $t > T$.
\end{theorem}

\begin{theorem}[Dimension selection]
$d = 3$ is the unique spatial dimension with both stable Keplerian orbits ($2 \leq d \leq 3$) and clean wave propagation (Huygens' principle: odd $d \geq 3$). Spacetime dimension $= 3 + 1 = 4$.
\end{theorem}

\textbf{Remark.} The physical argument dates to Ehrenfest (1917) and was elaborated by Tegmark (1997)~\cite{ehrenfest,tegmark}. What is new here is: (a)~the formal machine-checked proof (via \texttt{omega} on the arithmetic conditions), and (b)~the integration into the parent framework as a derived consequence rather than an independent observation.

\section{Concrete Computational Realizations}
\label{sec:realizations}

The following observables are reproduced in deterministic computational models. These results are empirical validations of the framework's concrete realizations, not formal Lean theorems.

\subsection{Gravitational observables}
\label{sec:gr}

\begin{table}[h]
\centering
\begin{tabular}{lcc}
\toprule
Observable & Measured & Expected (GR) \\
\midrule
Inverse-square exponent $\alpha$ & $2.004 \pm 0.063$ & $2.0$ \\
Raychaudhuri focusing & $\theta \propto b^{-0.983}$ & $b^{-1}$ \\
Shapiro time delay slope & $3.94$ & $4.0$ \\
Deflection angle & $\delta\phi \propto b^{-1.030}$ & $b^{-1}$ \\
Nonlinear amplification & $>10^6\times$ at strong $Q$ & Present \\
Horizon-like trapping & $b_{\text{crit}}$ grows with $Q$ & Present \\
\bottomrule
\end{tabular}
\caption{Gravitational observables from the $3{+}1$ causal diamond model.}
\label{tab:gr}
\end{table}

\subsection{Quantum benchmarks}
\label{sec:quantum_bench}

\begin{table}[h]
\centering
\begin{tabular}{lcc}
\toprule
Benchmark & Key result & Accuracy \\
\midrule
Double-slit fringe formula & $d\sin\theta = n\lambda$ & $0.1\%$ \\
Aharonov--Bohm & Shift $\propto$ flux, $2\pi$-periodic & $1.5\%$ \\
Bell/CHSH & $|S| = 2\sqrt{2}$ (Tsirelson bound) & Exact \\
Mach--Zehnder / eraser & Visibility 1.0/0.0/1.0 & Exact \\
GHZ/Mermin $N$-party & $|M| = 2^{N-1}$ & Exact \\
\bottomrule
\end{tabular}
\caption{Quantum benchmarks from geometric path amplitudes (computational, not Lean proofs). The phase prescription $e^{ikL}$ is derived from the source functional's linearity (Theorem~\ref{thm:propagation}); the wavenumber $k$ is the sole free parameter per particle type.}
\label{tab:quantum}
\end{table}

\section{Scaling Relation: The Cosmological Constant}

In the discrete causal-substrate picture, the cosmological constant is not a fixed Planck-scale vacuum energy density that must be fine-tuned away. It is a finite-number fluctuation of a causal counting measure.

For a Poisson causal set with density $\rho$ and causal past 4-volume $V$, the number of causal elements is $N = \rho V$. The Poisson fluctuation gives:
\[
\Lambda \sim \frac{1}{\sqrt{N}} = \frac{1}{\sqrt{\rho V}}, \qquad \Lambda^2 \cdot N = 1.
\]

For Planck density $\rho \sim l_P^{-4}$ and Hubble 4-volume $V_H$, $N \sim 10^{240}$, giving $\Lambda \sim 10^{-120}$ in Planck units---the correct order of magnitude for the observed cosmological constant. Standard QFT predicts $\Lambda \sim 1$ (wrong by $10^{120}$).

This is a scaling relation suggested by the discrete causal counting model, not yet a rigorous prediction. A full derivation requires tighter normalization and a precise definition of $V$ (causal past, Hubble, or horizon volume). The self-consistency equation (causal past volume depends on $\Lambda$) fixes the volume constant $c = 1/\rho$, introducing no additional free parameter.

A potentially testable consequence: $\Lambda$ is not exactly constant but fluctuates with $\delta\Lambda / \Lambda \sim \Lambda / 2$, distinguishable in principle from exactly constant dark energy.

\section{Derivation Integrity}
\label{sec:integrity}

We distinguish three categories:

\textbf{Formally proven (Lean 4, zero sorry).} Renormalization rigidity, null-cone determination ($1{+}1$ and general $n{+}1$ for $n \geq 2$, both formalized), null-link equivalence (algebraic, under density/boundedness conditions; full measure-theoretic Poisson derivation not formalized), Poisson counting uniqueness, source/dressing split (derived from source functional, with proven universal property), contracted Bianchi identity, Riemann tensor from metric data, complete 4D Lovelock uniqueness (contraction classification, Gauss--Bonnet vanishing, $\varepsilon$-exclusion, higher Lovelock vanishing), full nonabelian Yang--Mills equation and nonabelian Bianchi identity (Jacobi cancellation fully machine-checked), linear defect charge algebra, multiple conserved charges, spin-statistics structure, representation covariance, Born-type observable uniqueness, complex amplitude uniqueness (Frobenius), operational quantum hypotheses (each proven necessary), interference formula, Fourier decomposition of interference, discrete decoherence, Lindblad exponential decoherence with classical limit ($\varepsilon$-$\delta$ proven), dimension selection ($d = 3$ unique), primitive necessity (all 5 proven), dynamical stability (ker($L$) is SubspaceStable), variational Einstein (action density, non-degeneracy), Poisson substrate bridge, normalization theorem (1 free parameter $\rho$), cosmological constant scaling law ($\Lambda^2 N = 1$), propagation rule derived from source functional (linearity $\to$ additivity $\to$ multiplicative exponential $\to$ $e^{ikL}$, existence direction; uniqueness of characters not yet formalized), two-path interference formula $|z_1+z_2|^2 = 2 + 2\cos(k(s_1-s_2))$ (derived from propagation rule), physical selection principle (physical $\Leftrightarrow$ stationary, flat metric is a solution), chirality derived from K/P split plus gauge invariance of the source functional (\S\ref{sec:chirality}), anomaly cancellation uniqueness of Standard Model hypercharges (cubic factorization $18y_Q(y_d+4y_Q)(y_d-2y_Q)=0$, both solutions are the SM up to $u{\leftrightarrow}d$; \S\ref{sec:anomaly}), $\theta$-term parity-odd and orthogonal to parity-even source functional (\S\ref{sec:strongcp}), CP violation requires $N_g \geq 3$ generations (minimum for nonzero CKM phase; \S\ref{sec:generations}).

\textbf{Computationally verified (Python, deterministic).} GR observables in $3{+}1$ dimensions (Table~\ref{tab:gr}), quantum benchmarks (Table~\ref{tab:quantum}).

\textbf{Honestly primitive.} The source functional $\varphi$ is a framework primitive---the development derives what follows if one exists, not why nature should have one (though removing it gives a provably degenerate theory). The operational quantum hypotheses (quadratic regularity, rotation invariance, faithfulness, composition) are structural conditions, not derived from deeper operational axioms in the Hardy/CDP sense. The one continuous free parameter is $\rho$ (discreteness density). The Lie algebra (gauge group) is a structural choice; compactness (negative-definite Killing form) is derived from energy boundedness, but no bound on dim($\mathfrak{g}$) from the spacetime dimension is derived. The representation content of the fermion sector (which representations of the gauge group the matter fields transform under) is an input, not derived.

\textbf{Gauge structure: what is derived vs.\ assumed.} The framework naturally produces a \emph{two-origin} gauge structure: nonabelian factors arise from the Lie algebra primitive (constrained to be semisimple by Killing form non-degeneracy), while a U(1) factor arises independently from the source/dressing decomposition (the phase rotation $z \mapsto e^{i\theta}z$ is a consequence of the complex structure that Frobenius forces). The full gauge group is thus predicted to be reductive --- semisimple $\times$ U(1) --- matching the structure of the Standard Model group SU(3)$\times$SU(2)$\times$U(1)$_Y$, where the abelian hypercharge factor has a different geometric origin than the nonabelian color and isospin factors. The semisimple part is not selected by the framework. However, \emph{given} the representation content of one Standard Model generation (quarks as triplet-doublets, leptons as singlet-doublets, with the standard chirality assignment), anomaly cancellation uniquely determines all hypercharge ratios up to an overall normalization (Theorem~\ref{thm:anomaly-unique}).

\subsection{Chirality from the source/dressing decomposition}
\label{sec:chirality}

The source functional $\varphi$ is gauge-invariant (it is a framework primitive, defined before gauge structure is introduced). The K/P decomposition gives $\varphi(K(v)) = \varphi(v)$ (bridge) and $\varphi(P(v)) = 0$ (neutrality).  For any gauge transformation $g$ that preserves the K/P split, gauge invariance of $\varphi$ requires $\varphi(g_K(K(v))) = \varphi(K(v))$ --- the gauge action on the K-sector must preserve $\varphi$-values.  No such constraint applies to the P-sector, since $\varphi$ vanishes there regardless.

This asymmetry is chirality: the gauge group is \emph{more constrained} on the source sector (``left-handed'') than on the dressing sector (``right-handed'').  Any gauge transformation that acts nontrivially on the $\varphi$-charge must mix K and P components and therefore cannot preserve the split (\texttt{nontrivial\_gauge\_requires\_mixing} in \texttt{ChiralityFromKP.lean}).

Chirality is thus a \emph{consequence} of having a source functional (needed for gravity) and a nontrivial K/P split (needed for quantum interference), not an independent postulate.

\subsection{Anomaly cancellation and hypercharge uniqueness}
\label{sec:anomaly}

The defect sector produces charges via linear functionals (proven) and the K/P split provides the chirality grading derived above. Gauge anomaly cancellation --- the vanishing of $A_{abc} = \sum_i \chi_i \operatorname{Tr}[T_a^{(i)} \{T_b^{(i)}, T_c^{(i)}\}]$ --- is a consistency condition on the gauge--matter system, not an additional axiom.

\begin{theorem}[Anomaly uniqueness of SM hypercharges]
\label{thm:anomaly-unique}
Let $G = \mathrm{SU}(3) \times \mathrm{SU}(2) \times \mathrm{U}(1)_Y$ with one generation of fermions in the standard representations: quark doublets $Q_L$ (triplet-doublet), up-type singlets $u_R$ (triplet-singlet), down-type singlets $d_R$ (triplet-singlet), lepton doublets $L_L$ (singlet-doublet), and charged lepton singlets $e_R$ (singlet-singlet).  Denote the hypercharges of the corresponding left-handed Weyl fields by $(y_Q, y_u, y_d, y_L, y_e)$.

The four anomaly cancellation conditions --- $\mathrm{SU}(2)^2 \times \mathrm{U}(1)$, $\mathrm{SU}(3)^2 \times \mathrm{U}(1)$, gravitational, and cubic $\mathrm{U}(1)^3$ --- are:
\begin{align}
3y_Q + y_L &= 0, \label{eq:su2mixed}\\
2y_Q + y_u + y_d &= 0, \label{eq:su3mixed}\\
6y_Q + 3y_u + 3y_d + 2y_L + y_e &= 0, \label{eq:grav}\\
6y_Q^3 + 3y_u^3 + 3y_d^3 + 2y_L^3 + y_e^3 &= 0. \label{eq:cubic}
\end{align}
Conditions \eqref{eq:su2mixed}--\eqref{eq:grav} determine $y_L = -3y_Q$, $y_e = 6y_Q$, and $y_u + y_d = -2y_Q$.  Substituting into \eqref{eq:cubic} and factoring:
\[
18\, y_Q\, (y_d + 4y_Q)(y_d - 2y_Q) = 0.
\]
For $y_Q \neq 0$, the two solutions $y_d = 2y_Q$ and $y_d = -4y_Q$ both give the Standard Model hypercharge assignment (related by $u \leftrightarrow d$ relabeling):
\[
(y_Q,\, y_u,\, y_d,\, y_L,\, y_e) \;=\; y_Q \cdot (1,\, {-4},\, 2,\, {-3},\, 6).
\]
With the normalization $y_Q = 1/6$, this recovers $(1/6,\, {-2/3},\, 1/3,\, {-1/2},\, 1)$.
\end{theorem}

\noindent\textit{Lean verification.} Each step is machine-checked in \texttt{AnomalyConstraints.lean}: the four anomaly conditions for the SM charges (\texttt{sm\_all\_anomalies}), the linear reduction (\texttt{three\_conditions\_leave\_one\_dof}), the cubic factorization (\texttt{substituted\_cubic\_factors}, \texttt{quadratic\_factors}), the uniqueness (\texttt{anomaly\_uniqueness}), and the verification that both solutions satisfy all conditions (\texttt{both\_solutions\_are\_sm}).  Zero sorry, zero custom axioms.

\medskip\noindent The representation content --- that quarks are triplets and leptons are singlets of SU(3), that left-handed fields form doublets of SU(2) --- is an \emph{input}, not derived from the framework.  The theorem shows that this input, combined with anomaly cancellation, leaves zero freedom in the hypercharge assignments.

\medskip\noindent\textit{Enumeration of anomaly-free theories.} Among all chiral fermion sets using only fundamental representations of SU(3)$\times$SU(2) with arbitrary U(1) charges (exhaustively enumerated up to 20 Weyl fermions): 28 total anomaly-free structures exist; 24 use higher representations (adjoints, symmetric tensors); 4 use only fundamentals.  Of the fundamental-only structures, two use the SU(2) adjoint (triplet) for leptons.  With strict fundamentals (SU(3) triplet, SU(2) doublet), the Standard Model generation is the \emph{unique} anomaly-free structure, up to SU(3) conjugation.

The SM is also distinguished by a structural property: it is the minimal chiral theory where the cubic anomaly condition is \emph{algebraically independent} of the three linear conditions (rather than being automatically satisfied).  In all smaller chiral theories, the linear conditions alone determine the charge ratios and the cubic is redundant.  In the SM, the 2-dimensional null space of the linear conditions is cut to discrete branches by the cubic --- which is what produces the factorization of Theorem~\ref{thm:anomaly-unique}.

\medskip\noindent\textit{The conditional derivation chain.}  The full chain from framework to SM charge structure requires two non-derived inputs beyond the primitives: (1)~matter consists of elementary defects, i.e., the simplest localized excitations on the causal set, which transform in fundamental representations (a physical hypothesis motivated by the discrete substrate); (2)~a selection principle equivalent to requiring all anomaly conditions to be independently constraining (not derived, but consistent with the framework's general theme of maximal rigidity --- cf.\ Lovelock uniqueness, where every gravitational term independently constrains in 4D).  Given these two inputs plus the derived chirality (\S\ref{sec:chirality}), the SM generation is uniquely selected and all hypercharges are determined.

\subsection{Strong CP and the parity structure}
\label{sec:strongcp}

The QCD $\theta$-term $\mathcal{L}_\theta = (\theta/32\pi^2)\operatorname{Tr}[F_{\mu\nu}\tilde{F}^{\mu\nu}]$ involves the dual field strength $\tilde{F}^{\mu\nu} = \tfrac{1}{2}\varepsilon^{\mu\nu\rho\sigma}F_{\rho\sigma}$.  Since $\varepsilon$ is a pseudotensor, $\operatorname{Tr}[F\tilde{F}]$ is parity-odd: under the parity transformation $\varepsilon \to -\varepsilon$, it changes sign (\texttt{topological\_density\_parity\_odd}).  The Yang--Mills density $\operatorname{Tr}[F^2]$ is parity-even (no $\varepsilon$ dependence).

The source functional $\varphi$ is parity-even (it is derived from the trace of a symmetric matrix, which is invariant under spatial reflection).  A parity-odd quantity is orthogonal to any parity-even functional: if $q = -q$ under parity, then $q = 0$ under parity averaging.  Therefore, the $\theta$-term does not enter the source-sector field equations.

This provides a structural resolution of the strong CP problem: the $\theta$-parameter is not forbidden but is \emph{invisible} to the source functional that determines gravitational and gauge dynamics.  The resolution requires no axion, no new symmetry, and no fine-tuning --- it follows from the parity structure of the K/P decomposition, which is already present for independent reasons (quantum interference).

\medskip\noindent\textit{Remark.}  The discrete causal-set substrate plausibly lacks the smooth topology required for instanton sectors (which are classified by $\pi_3(G)$ and require smooth gauge field configurations).  If so, the $\theta$-parameter is not merely invisible but absent.  Verifying this requires computing the homology of Poisson causal sets using nerve complexes of Alexandrov neighborhoods, which remains open --- three simpler constructions (order complex, link complex, Vietoris--Rips) fail to recover manifold topology on $T^2$ validation tests.

\subsection{Generations and CP violation}
\label{sec:generations}

The quark mixing matrix for $N_g$ generations is $N_g \times N_g$ unitary, with $(N_g{-}1)(N_g{-}2)/2$ CP-violating phases.  For $N_g < 3$: zero phases, hence no CP violation.  For $N_g = 3$: exactly one phase.  CP violation is observed experimentally and is required for baryogenesis (Sakharov's second condition).  Therefore $N_g \geq 3$ (\texttt{cp\_violation\_requires\_d\_ge\_3}).

The framework's dimension selection theorem gives $d = 3$ as the unique spatial dimension permitting both stable orbits and clean (Huygens) wave propagation.  This is also the minimum dimension supporting CP violation: the same structural argument that selects $d = 3$ for gravitational stability gives the lower bound $N_g \geq 3$ for matter-antimatter asymmetry.  If $N_g$ tracks $d$ through any mechanism (e.g., the $d$ eigenvalues of the spatial metric perturbation providing $d$ independent defect propagation modes), then $N_g = 3$ follows from the dimension selection theorem.  This connection is conjectural; the exact value $N_g = 3$ is not derived.

\section{Discussion}

The framework presented here is one of the first formally verified developments attempting a common parent-structure treatment of gravity, gauge theory, matter, and quantum-style amplitude structure. The formal verification means the logical gap between hypotheses and conclusions is zero for the algebraic core.

Within the framework, the passage from classical to quantum-style behavior is not separately postulated but conditionally induced: if the parent object has a source functional (needed for a gravitational source term), and the dressing sector is nontrivial, then defects carry complex amplitudes, the Born-type observable is the unique rotation-invariant quadratic observable, and classical behavior arises from Lindblad-type dephasing.

\subsection{Limitations}

The null-link equivalence is established at the algebraic level under density and boundedness conditions; a full measure-theoretic derivation from Poisson processes is not formalized. The framework identifies the admissible tensor structure and a physical selection principle (stationary $\Leftrightarrow$ $G = 0$), but does not yet formalize the full space of solutions. The cosmological constant result is a scaling relation, not yet a rigorous prediction. The propagation rule derivation uses the existence direction of the characters-of-$\mathbb{R}$ theorem (exponential satisfies the conditions); the uniqueness direction (every continuous character is exponential) is classical and will be formalized in a subsequent version. No bound on gauge algebra dimension is derived from the spacetime dimension; the gauge group is a structural choice constrained only by compactness (energy boundedness) and non-degeneracy of the Killing form. The representation content of the fermion sector is an input; anomaly cancellation then uniquely determines the hypercharge assignment but does not select the gauge group itself.

\subsection{Open questions}

Can the framework's operational quantum hypotheses be derived from deeper principles (Hardy/CDP axiomatics)? Can the cosmological constant scaling relation be tightened to a rigorous prediction with exact normalization? Does the defect sector's structure force matter into fundamental representations, or is this an independent physical hypothesis?  Can the homology of Poisson causal sets be computed using nerve complexes of Alexandrov neighborhoods (three simpler constructions --- order complex, link complex, Vietoris--Rips --- have been computationally shown to fail on $T^2$)?  If discrete topology is computable, does the discrete second Chern class vanish, confirming the absence of instanton sectors? Does the number of fermion generations equal the spatial dimension ($N_g = d = 3$) through the eigenvalue structure of the spatial metric perturbation?  What determines the mass spectrum and coupling constants?

\section{Conclusion}

We have presented a formally verified parent-structure framework in which the gravity endpoint, gauge structure, matter algebra, and a quantum-amplitude layer are selected or conditionally induced from a small parent architecture, with explicit primitive assumptions and concrete realizations. The classification results show that within the stated classes, the Einstein$+\Lambda$ field equation, the nonabelian Yang--Mills equation, the charge algebra, and the Born-type observable are each uniquely determined. Each of five framework primitives is proven necessary. The propagation rule $e^{ikL}$ is derived from the source functional's linearity, chirality is derived from the K/P decomposition, and the $\theta$-term is shown to be parity-orthogonal to the source functional.  Given the Standard Model representation content as an explicit hypothesis, anomaly cancellation uniquely determines all hypercharge ratios, with the cubic condition factoring to give exactly the Standard Model (Theorem~\ref{thm:anomaly-unique}).  The full Standard Model charge structure follows from two non-derived inputs (elementary defects and a selection principle on anomaly independence) plus the framework's derived chirality.  A discrete causal-counting model suggests a scaling relation for the cosmological constant at the correct order of magnitude.

\begin{thebibliography}{99}
\bibitem{lean4} L.~de Moura and S.~Ullrich, ``The Lean 4 theorem prover and programming language,'' in \emph{Proc. CADE-28}, LNCS vol.~12699, Springer, 2021.
\bibitem{sorkin} R.~D.~Sorkin, ``Spacetime and causal sets,'' in \emph{Relativity and Gravitation: Classical and Quantum}, World Scientific, 1991.
\bibitem{lovelock} D.~Lovelock, ``The Einstein tensor and its generalizations,'' \emph{J.\ Math.\ Phys.}\ \textbf{12}(3):498--501, 1971.
\bibitem{ehrenfest} P.~Ehrenfest, ``In what way does it become manifest in the fundamental laws of physics that space has three dimensions?'' \emph{Proc.\ Amsterdam Acad.}\ \textbf{20}:200--209, 1917.
\bibitem{tegmark} M.~Tegmark, ``On the dimensionality of spacetime,'' \emph{Class.\ Quantum Grav.}\ \textbf{14}:L69--L75, 1997.
\end{thebibliography}

\end{document}
